The timing of the pricing problem determines when a feedback loop actually arises. In the one-period zero-coupon case, fixing face and the payoff mapping fixes the default event before the issue price is computed. The premium rule can therefore be evaluated after default probabilities are known. The haircut comparisons exploit exactly this recursive structure. A lower additional discount increases proceeds at the given face; under the common additive rule it multiplies every candidate face by the same factor, which leaves the maximizing face unchanged.

With long maturity or interest obligations linked to financing conditions, today's price can instead affect future service and refinancing. Purchases or collateral demand can change yields, those yields can change the promised-service schedule, and that schedule can alter future default risk. A simultaneous pricing and policy problem then requires a specification of who holds the claims, how outstanding coupons are set, and how the safe rate responds. Research on bond supply and duration risk, such as \textcite{GreenwoodVayanos2014}, motivates considering such channels; it does not identify the collateral elasticity $\xi$ or an instrument-specific eligibility rule in this paper.

Institutional interpretation should remain separate from the consolidated accounting. Debt managers choose contractual terms and maturity, while a central bank applies collateral and risk-management policies under its own mandate. An observed reduction in default exposure can be relevant to those decisions without implying joint maximization by a single authority. Conversely, providing liquidity against an eligible asset does not by itself remove its underlying climate exposure. The distinction between liquidity support and a change in the contractual payment schedule is why the main instrument comparison and the collateral assumptions are reported separately.
